% GENERATED FILE. Edit canonical lessons, then run npm run generate:books.
% canonical-source-hash: fnv1a64:43b983c64f18ea2b
% canonical-lessons: ML-C432-gavesanam, ML-C432-gavesakan, ML-C432-prabandham, ML-C432-prophasar, ML-C432-prabhasanam

\chapter{Research}
\label{ch:ml-research}
\hlchaptermodality{432}

\begin{chapteropening}
\textbf{By the end of this chapter:} \emph{I can say research, a researcher, a thesis, a professor and a lecture.}

Complete the last lesson of chapter 432: I can say research, a researcher, a thesis, a professor and a lecture.
\end{chapteropening}

\section[\texorpdfstring{\ml{ഗവേഷണം} (gavēṣaṇaṁ)}{ഗവേഷണം (gavēṣaṇaṁ)}]{\ml{ഗവേഷണം} (gavēṣaṇaṁ) — research}

\label{lesson:ML-C432-gavesanam}



\begin{quote}
Before the new one: say the Malayalam for to make headway, then the Malayalam for a system.
\end{quote}

\subsection*{You'll want to know: \ml{ഗവേഷണം}}
\textbf{\ml{ഗവേഷണം}} — \emph{gavēṣaṇaṁ} — \textquotedblleft{}research\textquotedblright{}.

\begin{quote}
\ml{അവൾ} \ml{കടലിലെ} \ml{മീനുകളെക്കുറിച്ച്} \ml{ഗവേഷണം} \ml{നടത്തുന്നു}.
\end{quote}

\begin{quote}
\emph{avaḷ kaṭalile mīnukaḷekkuṟiccŭ gavēṣaṇaṁ naṭattunnu.}
\end{quote}

\begin{quote}
She is doing research on the fish of the sea.
\end{quote}

To do research is \ml{ഗവേഷണം} \ml{നടത്തുക}.

\begin{cousinweb}
\textbf{\ml{ഗവേഷണം}} (\emph{gavēṣaṇaṁ}) is borrowed from Sanskrit \emph{gaveṣaṇa}, \textquotedblleft{}desiring ardently\textquotedblright{}, the noun of \emph{gaveṣ}, \textquotedblleft{}to seek, to search for\textquotedblright{} (Monier-Williams). Monier-Williams files the verb under \emph{go}, \textquotedblleft{}a cow\textquotedblright{}, which is often read as a search for cattle; the research sense is modern.
\end{cousinweb}

\begin{grammarlens}[title={What you've built}]
One more word for giving reasons and opinions.
\end{grammarlens}

\subsection*{Your turn}
\begin{itemize}
\raggedright
  \item \emph{Say it:} \emph{gavēṣaṇaṁ}
  \item \emph{Say it:} \emph{gavēṣaṇaṁ}, once more
  \item \emph{Say it:} \emph{avaḷ kaṭalile mīnukaḷekkuṟiccŭ gavēṣaṇaṁ naṭattunnu.} — She is doing research on the fish of the sea.
  \item \emph{Recall:} say the Malayalam for to make headway, then the Malayalam for a system, then say \emph{gavēṣaṇaṁ} again
\end{itemize}

\begin{tcolorbox}[breakable,colback=teal!4,colframe=teal!35!black,title={Before you move on}]
What does \ml{ഗവേഷണം} mean? (\textquotedblleft{}Research\textquotedblright{}.) Say it once more.
\end{tcolorbox}
\section[\texorpdfstring{\ml{ഗവേഷകൻ} (gavēṣakan)}{ഗവേഷകൻ (gavēṣakan)}]{\ml{ഗവേഷകൻ} (gavēṣakan) — a researcher}

\label{lesson:ML-C432-gavesakan}



\begin{quote}
Before the new one: say the Malayalam for a system, then the Malayalam for research.
\end{quote}

\subsection*{You'll want to know: \ml{ഗവേഷകൻ}}
\textbf{\ml{ഗവേഷകൻ}} — \emph{gavēṣakan} — \textquotedblleft{}a researcher\textquotedblright{}.

\begin{quote}
\ml{ആ} \ml{ഗവേഷകൻ} \ml{പത്ത്} \ml{വർഷം} \ml{കാട്ടിൽ} \ml{ജീവിച്ചു}.
\end{quote}

\begin{quote}
\emph{ā gavēṣakan pattŭ varṣaṁ kāṭṭil jīviccu.}
\end{quote}

\begin{quote}
That researcher lived in the forest for ten years.
\end{quote}

Plural: \ml{ഗവേഷകർ}, \emph{gavēṣakar}.

Plural: \ml{ഗവേഷകർ}.

\begin{cousinweb}
\textbf{\ml{ഗവേഷകൻ}} (\emph{gavēṣakan}) is built on the same Sanskrit word as \ml{ഗവേഷണം}, \emph{gaveṣaṇa}, \textquotedblleft{}searching\textquotedblright{} (Monier-Williams), with the agent ending -\ml{കൻ}, \textquotedblleft{}one who\textquotedblright{}.
\end{cousinweb}

\begin{grammarlens}[title={What you've built}]
One more word for giving reasons and opinions.
\end{grammarlens}

\subsection*{Your turn}
\begin{itemize}
\raggedright
  \item \emph{Say it:} \emph{gavēṣakan}
  \item \emph{Say it:} \emph{gavēṣakan}, once more
  \item \emph{Say it:} \emph{ā gavēṣakan pattŭ varṣaṁ kāṭṭil jīviccu.} — That researcher lived in the forest for ten years.
  \item \emph{Recall:} say the Malayalam for a system, then the Malayalam for research, then say \emph{gavēṣakan} again
\end{itemize}

\begin{tcolorbox}[breakable,colback=teal!4,colframe=teal!35!black,title={Before you move on}]
What does \ml{ഗവേഷകൻ} mean? (\textquotedblleft{}A researcher\textquotedblright{}.) Say it once more.
\end{tcolorbox}
\section[\texorpdfstring{\ml{പ്രബന്ധം} (prabandhaṁ)}{പ്രബന്ധം (prabandhaṁ)}]{\ml{പ്രബന്ധം} (prabandhaṁ) — a thesis; a paper}

\label{lesson:ML-C432-prabandham}



\begin{quote}
Before the new one: say the Malayalam for research, then the Malayalam for a researcher.
\end{quote}

\subsection*{You'll want to know: \ml{പ്രബന്ധം}}
\textbf{\ml{പ്രബന്ധം}} — \emph{prabandhaṁ} — \textquotedblleft{}a thesis; a paper\textquotedblright{}.

\begin{quote}
\ml{അവൻ} \ml{തന്റെ} \ml{പ്രബന്ധം} \ml{സർവകലാശാലയിൽ} \ml{സമർപ്പിച്ചു}.
\end{quote}

\begin{quote}
\emph{avan tanṟe prabandhaṁ sarvakalāśālayil samarppiccu.}
\end{quote}

\begin{quote}
He submitted his thesis to the university.
\end{quote}

A research paper read at a meeting is also a \ml{പ്രബന്ധം}.

\begin{cousinweb}
\textbf{\ml{പ്രബന്ധം}} (\emph{prabandhaṁ}) is borrowed from Sanskrit \emph{prabandha}, \textquotedblleft{}a band, a connexion\textquotedblright{} (CDIAL 8695), and in later Sanskrit a connected composition (Monier-Williams).
\end{cousinweb}

\begin{grammarlens}[title={What you've built}]
One more word for giving reasons and opinions.
\end{grammarlens}

\subsection*{Your turn}
\begin{itemize}
\raggedright
  \item \emph{Say it:} \emph{prabandhaṁ}
  \item \emph{Say it:} \emph{prabandhaṁ}, once more
  \item \emph{Say it:} \emph{avan tanṟe prabandhaṁ sarvakalāśālayil samarppiccu.} — He submitted his thesis to the university.
  \item \emph{Recall:} say the Malayalam for research, then the Malayalam for a researcher, then say \emph{prabandhaṁ} again
\end{itemize}

\begin{tcolorbox}[breakable,colback=teal!4,colframe=teal!35!black,title={Before you move on}]
What does \ml{പ്രബന്ധം} mean? (\textquotedblleft{}A thesis; a paper\textquotedblright{}.) Say it once more.
\end{tcolorbox}
\section[\texorpdfstring{\ml{പ്രൊഫസർ} (prophasar)}{പ്രൊഫസർ (prophasar)}]{\ml{പ്രൊഫസർ} (prophasar) — a professor}

\label{lesson:ML-C432-prophasar}



\begin{quote}
Before the new one: say the Malayalam for a researcher, then the Malayalam for a thesis.
\end{quote}

\subsection*{You'll want to know: \ml{പ്രൊഫസർ}}
\textbf{\ml{പ്രൊഫസർ}} — \emph{prophasar} — \textquotedblleft{}a professor\textquotedblright{}.

\begin{quote}
\ml{പ്രൊഫസർ} \ml{ക്ലാസിൽ} \ml{ഒരു} \ml{കഥ} \ml{പറഞ്ഞു}.
\end{quote}

\begin{quote}
\emph{prophasar klāsil oru katha paṟaññu.}
\end{quote}

\begin{quote}
The professor told a story in class.
\end{quote}

Plural: \ml{പ്രൊഫസർമാർ}.

\begin{cousinweb}
\textbf{\ml{പ്രൊഫസർ}} (\emph{prophasar}) was borrowed from English \emph{professor}.
\end{cousinweb}

\begin{grammarlens}[title={What you've built}]
One more word for giving reasons and opinions.
\end{grammarlens}

\subsection*{Your turn}
\begin{itemize}
\raggedright
  \item \emph{Say it:} \emph{prophasar}
  \item \emph{Say it:} \emph{prophasar}, once more
  \item \emph{Say it:} \emph{prophasar klāsil oru katha paṟaññu.} — The professor told a story in class.
  \item \emph{Recall:} say the Malayalam for a researcher, then the Malayalam for a thesis, then say \emph{prophasar} again
\end{itemize}

\begin{tcolorbox}[breakable,colback=teal!4,colframe=teal!35!black,title={Before you move on}]
What does \ml{പ്രൊഫസർ} mean? (\textquotedblleft{}A professor\textquotedblright{}.) Say it once more.
\end{tcolorbox}
\section[\texorpdfstring{\ml{പ്രഭാഷണം} (prabhāṣaṇaṁ)}{പ്രഭാഷണം (prabhāṣaṇaṁ)}]{\ml{പ്രഭാഷണം} (prabhāṣaṇaṁ) — a lecture; a speech}

\label{lesson:ML-C432-prabhasanam}



\begin{quote}
Before the new one: say the Malayalam for a thesis, then the Malayalam for a professor.
\end{quote}

\subsection*{You'll want to know: \ml{പ്രഭാഷണം}}
\textbf{\ml{പ്രഭാഷണം}} — \emph{prabhāṣaṇaṁ} — \textquotedblleft{}a lecture; a speech\textquotedblright{}.

\begin{quote}
\ml{അദ്ദേഹത്തിന്റെ} \ml{പ്രഭാഷണം} \ml{ഒരു} \ml{മണിക്കൂർ} \ml{നീണ്ടു}.
\end{quote}

\begin{quote}
\emph{addēhattinṟe prabhāṣaṇaṁ oru maṇikkūr nīṇṭu.}
\end{quote}

\begin{quote}
His lecture lasted an hour.
\end{quote}

The speaker is \ml{പ്രഭാഷകൻ}.

\begin{cousinweb}
\textbf{\ml{പ്രഭാഷണം}} (\emph{prabhāṣaṇaṁ}) is borrowed from Sanskrit \emph{prabhāṣaṇa}, \textquotedblleft{}explanation\textquotedblright{} (Monier-Williams), from \emph{pra-bhāṣ}, \textquotedblleft{}to speak out\textquotedblright{}.
\end{cousinweb}

\begin{grammarlens}[title={What you've built}]
One more word for giving reasons and opinions.
\end{grammarlens}

\subsection*{Your turn}
\begin{itemize}
\raggedright
  \item \emph{Say it:} \emph{prabhāṣaṇaṁ}
  \item \emph{Say it:} \emph{prabhāṣaṇaṁ}, once more
  \item \emph{Say it:} \emph{addēhattinṟe prabhāṣaṇaṁ oru maṇikkūr nīṇṭu.} — His lecture lasted an hour.
  \item \emph{Recall:} say the Malayalam for a thesis, then the Malayalam for a professor, then say \emph{prabhāṣaṇaṁ} again
\end{itemize}

\begin{tcolorbox}[breakable,colback=teal!4,colframe=teal!35!black,title={Before you move on}]
What does \ml{പ്രഭാഷണം} mean? (\textquotedblleft{}A lecture; a speech\textquotedblright{}.) Say it once more.
\end{tcolorbox}
